\section{Discussion and limitations}\label{sec:discussion}

The numerical results above are conditional on choices that the working
paper makes explicit; this section catalogues them.

\paragraph{Granular default data.}
The Visale and ANIL series on residential rent arrears are released as
annual aggregates rather than per-apartment time-stamps. We use them as a
prior for the marginal default probability $p$ but cannot calibrate the
copula correlation $\rho$ directly --- a single building yields a single
realisation of joint defaults, so $\rho$ remains essentially unidentified
from the building alone. We therefore quote results at the prior value
$\rho = \rhoBaseline{}$ and report the sensitivity to
$\rho \in \{0.00, 0.10, 0.20, 0.30, 0.45\}$ in
\texttt{notebooks/05\_stress\_and\_sensitivity.ipynb}.

\paragraph{Vasicek long-run mean.}
The MLE on French OAT 10Y delivers $\hat b = \paramVasicekB{}$ ---
substantially above the current short-rate level because the historical
sample (1960--2026) spans the 1970s/1980s high-rate regime. For long-run
discount factors this is a feature; for shorter horizons one would
typically work off the current rate curve. We document the implication in
the Sharpe ratios reported above: when the long-run discount factor is
high relative to the senior coupon, the tranche annualised return relative
to the risk-free benchmark is mechanically negative even though the
tranche prices close to par. This is the calibration trade-off, not a
defect of the structure.

\paragraph{Identifiability of jumps.}
Numerical MLE on the Merton (1976) jump-diffusion does not improve AIC on
the Paris sample (see \cref{tab:calibration}); the jump intensity and the
jump-size standard deviation collapse together, which is a well-known
identifiability issue when the data sample has too few clean jumps. We
keep GBM as the baseline and view the jump extension as a robustness
check.

\paragraph{Static vs dynamic correlation.}
The Gaussian / Student-t copulas embed default correlation through a
single asset-value draw at $t = 0$. The Cox doubly stochastic model
relaxes this by letting the hazard rate evolve over time, producing
time-varying default correlation. The numerical gap between the static
copulas and the Cox model (\cref{fig:loss-dist}) is small in the body
of the loss distribution but visible in the tail; the choice depends
on the application horizon.

\paragraph{Insurance pricing endogeneity.}
We treat the insurance premium as an exogenous cost subtracted from each
period's rent. In reality the insurer would re-underwrite over time and
the premium would adjust to realised losses. Our coverage-cap-times-loss
construction is an internal-consistency benchmark, not a market premium
quote.

\paragraph{No equilibrium argument.}
We do not attempt a full equilibrium price for the three vehicles; the
fair prices we report are physical-measure expected discounted cash flows
on the calibrated paths. A risk-neutral pricing would require either
traded comparable instruments (which do not exist for this contract
structure) or an explicit risk-aversion assumption.

\paragraph{Fair-coupon existence.}
The Brent solver of \cref{ssec:fair-coupons} returns a finite senior
fair coupon under the Gaussian and Student-t copulas but no solution
under the Cox model and no mezzanine solution under any of the three:
the available rent net of maintenance and senior coupon is not
sufficient to bring mezz to par within the considered coupon bracket.
We emit \texttt{NaN} in \texttt{artifacts/fair\_coupons.csv}. The
economic interpretation is that the cash-flow waterfall is too senior-
heavy given the calibrated $(g - m)$ spread; standard structural fixes
(over-collateralisation, cash trap, principal acceleration) would
restore mezz solvability and are kept out of scope for the current
working paper.

\paragraph{Fair-coupon vs current quoted spreads.}
Our fair-coupon solution treats the underlying obligor pool as
exogenously calibrated. In a market with traded comparable instruments
(e.g.\ tranches on iTraxx Europe) one would also calibrate $\rho$
implicitly from quoted spreads, the so-called ``base correlation''
approach \autocite{AndersenSideniusBasu2003}. We do not pursue this here
because no comparable contract trades on Paris residential rent.

\paragraph{Joint vs sequential fair-coupon solver.}
The headline pipeline runs both a sequential Brent solver (senior coupon
first, then mezzanine conditional on the senior fair coupon) and a fully
joint 2-D \texttt{fsolve} on $(c_S, c_M)$. On the calibrated Paris
baseline both solvers concur: a finite fair coupon exists for the senior
($\sim 3.8$\,\%) but \emph{no} feasible pair $(c_S, c_M)$ brings both
tranches to par, irrespective of the optimisation routine. The
sequential solver therefore reports a numerically valid senior fair
coupon and a \texttt{NaN} mezzanine; the joint solver fails to converge
and reports both as \texttt{NaN}. Either way the conclusion is the same:
the rental cash flow available after maintenance is structurally
insufficient to remunerate the mezzanine layer at par on a stand-alone
basis.

\paragraph{Effect of the overcollateralisation test.}
Activating the OC test (\texttt{waterfall.oc\_test\_enabled = true} in
the YAML) diverts the equity periodic cash flow to a reserve account
whenever the OC ratio falls below the trigger (1.10 by default) and
releases the reserve once the OC ratio recovers above the target (1.15).
On the Paris baseline the test sees infrequent triggers and the
quantitative impact on senior is modest ($<$ 10 bps on fair-to-par);
on the GFC overlay the trigger fires routinely and the cash trap
materially supports the mezz tranche at the expense of equity. The
mechanic is exactly the structural enhancement standard CDO offerings
include — our implementation makes the empirical impact measurable.

\paragraph{Cox calibration on real unemployment data.}
The Cox doubly-stochastic model now uses a CIR macro factor calibrated
on the INSEE French ILO unemployment series 1975Q1–2026Q1 ($n = 204$
observations). The fitted $(\hat\kappa, \hat\theta, \hat\xi)$ deliver a
long-run mean around 8.7\,\% — the historical sample mean — and satisfy
the Feller positivity condition. Compared to the static asset-value copulas, the data-driven Cox
calibration delivers a tighter senior fair-to-par (Cox $0.896$ vs
Gaussian $0.939$, a $\sim 4.3$ pp gap that exceeds the senior bootstrap
CI width by two orders of magnitude) and a much wider mezzanine
($0.262$ vs $0.786$). The mechanism is the macro factor's persistence:
time-varying default correlation creates loss clustering that the
static asset-value copulas miss, biting deepest at the mezzanine
attach point.

\paragraph{Gaussian vs Student-t senior coupon — quantified.}
Across the headline calibrations the senior fair coupon under the
Student-t copula sits about twelve basis points above the Gaussian.
With $\nu = 5$ and $\rho = \rhoBaseline{}$ the implied tail dependence
$\lambda_L \approx 0.08$ raises the joint senior-loss probability
relative to the Gaussian benchmark, whose tail dependence is zero by
construction.

\paragraph{Future work.}
The natural next steps are (i) a Hull-White / shifted-Vasicek extension
that admits a fitted initial yield curve, (ii) richer Beta-LGD
parameterisation calibrated on micro-level Visale claims data when
available, and (iii) a Bayesian sensitivity on $\rho$ with an informative
prior anchored on macro arrears statistics.
